% !TEX root = ../Report.tex
% ======================================================================
\chapter{Background and Related Work}
% ======================================================================

This chapter positions GiTrip within existing trip-planning and collaboration
tooling, and introduces foundational techniques used in the system: version
control semantics, scheduling heuristics, string matching, accessibility
heuristics, and relevant web standards.

\section{Trip Planning and Routing Tools}
Most users plan trips using mapping and routing tools such as Google Maps and
Apple Maps \cite{ref_google_maps,ref_apple_maps}. These applications provide
high-quality route computation and support multiple transport modes, but their
data model is fundamentally \textit{single-plan centric}. Users can create lists
and saved routes, but they cannot easily branch into alternative itineraries,
compare versions, and merge parallel edits with structured conflict handling.

Public-transport-focused tools such as Citymapper offer strong transit routing
and navigation \cite{ref_citymapper}, but similarly treat the itinerary as an
output rather than a version-controlled artefact. Trip planning products such as
Wanderlog, Furkot, Roadtrippers, and TripIt provide itinerary management
features and collaboration to varying degrees
\cite{ref_wanderlog,ref_furkot,ref_roadtrippers,ref_tripit}, yet they generally
emphasise itinerary presentation and booking integration rather than explicit
version control semantics. Google Travel aggregates trip components and
suggestions but does not provide Git-style change management
\cite{ref_google_travel}.

GiTrip's differentiation is therefore not \textit{better routing} than these
tools; instead, it provides structured change management and alternative
exploration, paired with sufficient routing and scheduling integration to keep
plans feasible.

\section{Generic Collaboration Tools}
When collaboration is required, users commonly use documents and spreadsheets to
manage itineraries: Google Docs/Sheets, Microsoft Word/Excel, and increasingly
general workspaces like Notion
\cite{ref_google_docs,ref_google_sheets,ref_ms_word,ref_ms_excel,ref_notion}.
These platforms have excellent collaboration features but treat an itinerary as
unstructured text or loosely structured tables. As a result, they cannot
automatically recompute travel times after edits, offer interactive map or
timeline visualisation, detect domain-specific conflicts such as overlapping
times, or provide branching semantics for exploring alternatives.

GiTrip's objective is to bring the strengths of version control
(branching/merging/history) into a domain-specific planner where plans are
first-class structured data.

\section{Version Control Concepts and Applicability}
Git is the dominant distributed version control system in software development
\cite{ref_git}. Its key semantics include commits (immutable versions), branches
(named pointers to commits), and merges (reconciling divergent commit histories).
Git's internal model stores snapshots rather than deltas: each commit represents
a complete view of the repository state \cite{ref_progit_snapshots}. This is
particularly relevant for GiTrip because a trip plan can be represented naturally
as a JSON document. Snapshot storage enables instant branch switching (render the
head snapshot), version restoring (create a new commit from an older snapshot),
and three-way merges (compare a base snapshot against two branch snapshots).

However, applying Git semantics to structured itinerary data introduces
challenges: traditional line-based diff/merge strategies (e.g., \texttt{diff3})
are designed for text and can produce poor results for structured JSON. The
fundamental three-way merge idea remains useful: derive changes in two branches
relative to a common ancestor and combine them, surfacing conflicts when both
sides change the same part incompatibly. Formal work on \texttt{diff3} analyses
its behaviour as a widely used three-way merge algorithm \cite{ref_diff3_paper}.

GiTrip therefore adapts the three-way merge concept to structured itinerary
objects rather than line-based text.

\section{Scheduling and Route Optimisation}
A key sub-problem in trip planning is ordering stops efficiently. This relates to
the Travelling Salesman Problem (TSP), which is computationally hard in general
and commonly addressed with heuristics
\cite{ref_tsp_study,ref_tsp_local_search}. GiTrip uses heuristics that provide
good-enough results for interactive planning rather than exact optimal solutions,
which aligns with practical timetabling and scheduling approaches discussed in
the automated timetabling literature \cite{ref_patat}.

For example, the Nearest Neighbour (NN) heuristic greedily selects the closest
next stop and is a common baseline; it is known to be simple and fast while not
always optimal \cite{ref_nn_tsp}. GiTrip also includes an educational algorithm
visualisation that demonstrates NN, NN + 2-opt improvement
\cite{ref_nn_2opt_hybrid}, and Held-Karp
dynamic programming for small instances, helping explain why the production
scheduler uses heuristics.

\section{Approximate String Matching for Place Search}
Place search relies on external geocoding and search services (e.g., Nominatim).
Returned results can be noisy or differently formatted across user inputs. GiTrip
improves usability by re-ranking candidates using approximate string matching,
normalising input strings and applying edit distance techniques
\cite{ref_navaro_string_matching}. This improves the ``easy add'' flow where
users paste free-form place names and expect robust matching.

\section{Structured Data Merge and JSON Patch Formats}
JSON-based systems frequently represent updates using patch formats. RFC 7396
(JSON Merge Patch) defines a merge-based approach for describing modifications
to JSON objects \cite{ref_rfc7396}. However, general-purpose patch formats do
not provide domain-specific semantics for itineraries: for example, they do not
encode constraints such as time overlaps or opening-hours validity. GiTrip
therefore uses a domain-aware merge strategy:
\begin{itemize}
  \item file-like maps (\texttt{snapshot.files}) are merged per-key using deep
        equality,
  \item trip plans (\texttt{snapshot.plan}) are merged per-day and per-stop with
        explicit conflict detection for semantic fields.
\end{itemize}

\section{Usability and Accessibility Frameworks}
GiTrip must balance power-user semantics (branches, merges) with novice
usability. Nielsen's usability engineering framework provides well-established
heuristics for evaluating interfaces and diagnosing common usability problems
\cite{ref_nielsen}. Nielsen's heuristics were preferred over heavier frameworks
because they are practical for small teams, support rapid iteration, and directly
map to UI design decisions such as visibility of system state, consistency, error
prevention, and recognition over recall.

Accessibility was guided by WCAG 2.2 principles \cite{ref_wcag22}, particularly
around keyboard navigation, contrast, and clear labelling. GiTrip's use of
interactive map and timeline components introduces accessibility challenges; this
report therefore discusses mitigations and remaining limitations.

For evaluation, GiTrip used a survey instrument that included the System
Usability Scale (SUS) and Likert-style questions aligned with constructs from the
Technology Acceptance Model (TAM), notably perceived usefulness and perceived
ease of use \cite{ref_tam_davis}. TAM was not used as a strict statistical model
due to limited sample size, but its constructs informed question design and
interpretation.

\section{Summary: Gap in Existing Approaches}
In summary, existing trip planning tools provide routing and itinerary management
but lack robust version control semantics. Existing collaboration tools provide
sharing but lack domain structure. Version control systems provide branching and
merging but do not operate on itinerary semantics. GiTrip's novelty is that it
combines these approaches into a coherent system: a version-controlled,
constraint-aware, user-editable itinerary planner.
